\scanpage{29}
\mspage{13}
correspondantes de relèvements \inserted{de $\{1,\tau\}$} sont resp. les relèvements en $\tau_0$, en $\tau_1$ et en $\tau_\infty$ --- qui ne donnent rien de nouveau. Reste à considérer $\tau\dot\sigma_1$ (et ses conjugués \inserted{par $\rho$} $\tau\dot\sigma_\infty$, $\tau\dot\sigma_0$), on a
\[
\dot\sigma_1(2)=1/2,\quad\text{donc}\qquad(\tau\dot\sigma_1)(2)=\frac12
\]
et
\[
U^{\tau\dot\sigma_1}=\{z\in U\mid z\bar z=1\}=\mathbb U-\{1\},
\]
où $\mathbb U=\{z\in\mathbb C\mid|z|=1\}$. Cet ensemble contient $P$ et $\bar P$, et ne donne donc pas de nouveaux relèvement p.r. à celui déjà donné, avec $\tilde\sigma_1$ relevant $\dot\sigma_1\tau$.

Il faudrait quand-$\unsure{\bar m|=}$ donner aussi une description de $\mathcal E$ tout entier par gén. et relations, en partant d'un syst. de gén. de $\Gamma=\unsure{\mathfrak S_3|\mathcal S_3}\times\mathbb Z/2$, qu'on va relever de façon judicieuse. On prendra les \inserted{quatre} générateurs
\deleted{
\[
\unsure{\rho'|\dot\rho'}=\rho\tau,\quad\deleted{\unsure{\tau}}\underset{\substack{\Vert\\\dot\sigma_0\tau}}{\widetilde{\dot\sigma}_0},\quad\underset{\substack{\Vert\\\dot\sigma_1\tau}}{\widetilde{\dot\sigma}_1},\quad\underset{\substack{\Vert\\\dot\sigma_\infty\tau}}{\widetilde{\dot\sigma}_\infty},
\]
\inserted{(donc ${\rho'}^6=1$, ${\rho'}^3=\tau$, ${\rho'}^2=\rho^2=\rho^{-1}$)}
avec les relations
\[
\left\{\begin{aligned}
\widetilde{\dot\sigma}_1^2&=\widetilde{\dot\sigma}_2^2=\widetilde{\dot\sigma}_3^2=1\\
\widetilde{\dot\sigma}_0\widetilde{\dot\sigma}_1&=\widetilde{\dot\sigma}_1\widetilde{\dot\sigma}_\infty=\widetilde{\dot\sigma}_\infty\widetilde{\dot\sigma}_0=\unsure{{\rho'}^{\deleted{-1}2}|\dot\rho^{-2}}\\
{\rho'}^6&=1\\
{\rho'}^3&\text{ commute aux }\dot\sigma_i
\end{aligned}\right.
\]
}
\mspage{28}
